
\section{Overview}
\label{ch:overview}
In this research we will consider a simplified model of a relative modular robot 
and corresponding environment, designed with a potential for self-replication 
and based on the design, and material-robot principle, 
developed by  \citeauthor{Jenett2019MaterialRobot} \cite{Jenett2019MaterialRobot} 
and  \citeauthor{AbdelRahman2022Hierarchical} \cite{AbdelRahman2022Hierarchical}.

In our model, as in the original work, robot (see Figure \ref{fig:VoxelAndRobotOriginal}b)
operates by manipulating and connecting parts (see Figure \ref{fig:VoxelAndRobotOriginal}a)
of the same type as its own building blocks. 
These parts can be considered as material for assembling different structures
and are distributed in the environment in ordered or arbitrary manner. 

Despite their functional diversity, all blocks and parts
are geometrically identical and take the form of cuboctahedral voxels,
general view of it is shown in Figure \ref{fig:Voxel}.

\begin{figure}[h!]
    \centering
    \includegraphics[width=7cm]{Voxel.png} 
    \caption{Material part or single robot block geometry in Gazebo simulation environment}
    \label{fig:Voxel}
\end{figure}

Further in text, we refer to elements attached to the robot as "blocks"
and call the same entities when distributed in the environment and 
not connected to the robot as "parts." 

The robot  is composed of a sequence of blocks of
different functional types, an example of a robot can be seen on 
Figure \ref{fig:RobotStructureGazebo}, where different types of blocks marked with different colors.
The first and the last blocks of the sequence serve as end-effectors equipped with 
grippers, enabling the robot to interact with its environment. The intermediate 
blocks constitute a rigid structure or act as joints, which may provide rotations 
around different axis, for example changing yaw or pitch angles. 
\begin{figure}[h!]
    \centering
    \includegraphics[width=11cm]{RobotStructureGazebo.png} 
    \caption{Robot structure in Gazebo simulation environment}
    \label{fig:RobotStructureGazebo}
\end{figure}

This uniform geometry facilitates a modular construction paradigm, 
where the same physical units can be rearranged to produce different structures and
robot morphologies.

At any given time, robot in our model is capable of handling a single part, 
aligning it with the existing discrete environment and attaching it to some structure 
or new robot. 
Through repeated execution of this process, the robot can in principle assemble a 
replica of itself, thereby demonstrating the fundamental process of self-replication. 

The robot operates within a structured environment defined as a lattice 
field that is geometrically compatible with the cuboctahedral voxel design 
this field represented as while base structure underneath the robot 
on a Figure \ref{fig:RobotStructureGazebo}. 
This lattice  constrains both the robot and the individual 
parts to occupy discrete positions, ensuring that all spatial 
relations, and connections can be described in terms of 
transitions between well-defined integer coordinates. 
Such discretization not only simplifies the modeling of robot kinematics 
and assembly operations but also provides a natural basis for formal verification, 
since the system state space becomes finite and directly available for 
exploration by model-checking techniques.

In this chapter, we focus on methods to formally analyze the behavior of 
such systems. Model checking provides a rigorous framework for verifying 
whether the system satisfies specific properties of interest. In particular, 
we investigate the use of symbolic model checker - NuSMV, to evaluate 
both safety and liveness conditions for the modular robot using linear temporal logic (LTL)
and brunching temporal logic (CTL). 
Safety properties ensure that undesirable states, such as collisions or 
inconsistent connections, are never reached. 
Liveness properties, on the other hand, guarantee that the system is always capable of making 
progress, for example, that a new block can eventually be attached to 
a target or that a new robot can continue self-replication process.

Beyond simple verification, we also explore the potential of using NuSMV
counterexample generation as a decision-making strategy. Counterexamples provided 
by the model checker correspond to concrete sequences of state transitions that 
lead to a violation of a previously negated property. 
Interpreted constructively, such sequences 
can serve as action plans or decision sequences for the robot, 
indicating feasible operations that allow it to achieve its objectives. 
This approach establishes a bridge between formal verification and automated planning
and can increase the reliability and effectiveness of the system, and also
add flexibility to decision-making in complex and dynamic conditions. 

% \textcolor{red}{ TODO: Add individual block geometry figure }

\section{Analysis of the system}

According to the classification given in Chapter \ref{ch:SRRSclassification} 
described self-replicating system for 
the general configuration of a robot can be classified as active system with 
hierarchical type of replication and capability for direct and indirect replications. 

It has hierarchical replication type because in general it can produce 
robots capable of self-replication as well as robots or mechanical structures 
lacking of such capability.

It is active since it is directly engaged in constructing child robots. 

The system can be homogeneous or heterogeneous based on initial conditions and programmed 
behavior. 
Homogeneous type of SRRS consist of robotically identical units, whereas heterogeneous 
enables generation of multiple robot types, with potential for specialization.
Therefore, our system will  be homogeneous if we will use only one type of robot with predefined 
structure that will replicate itself and heterogeneous if there will be multiple types 
on start or during replication.

In this chapter we will assume the system is homogeneous but will 
analyze effect of different types or configuration of robots in different environments
on the properties of a system.

Using model proposed by \citeauthor{Lee2007Robotic} in \cite{Lee2007Robotic} 
and described in Chapter \ref{ch:srrsModels}  self-replicating process in our system 
can be formally represented as equation \ref{eq:formalRepresentation}.
\begin{equation}
    \label{eq:formalRepresentation}
    (R' , M' , E ) = \Phi(R, M, E, t) 
\end{equation}

Where:\\
$M$ -- multiset (set with repetitions) of available parts (material) located in the environment
to be used for building replicas robots by an initial system.\\
$R$ --  a multiset of initial robots to be reproduced. In our case one robot with predefined 
configuration. Each element of $R$ includes specific robot structure in form of 
sequence of blocks, geometry, kinematic limitations and robot control system.\\
$E$ -- is a multiset of environmental structures and/or elements involved in the replication
process (but not replicated). In our experiment is represented by base and optional obstacles.\\
$R',M'$ -- the new robots (or build structures) and the remaining and altered material respectively.\\
$\Phi$ -- the replication function of time $t$.

This model is a special case of a Formula \ref{eq:formalModel2phi} for completely structured
environment. $E=E'$ since each element of the environment (base) is either fixed at a precise  location,
or described in integer coordinates and there will be  no modification or change 
in any structures during the self-replication process.

This formulation enables modeling dynamic replication scenarios over time and accounts 
for state changes in materials and  robots.

% To sum up the  description above we have a material-robot system, 
% where the environment material unit introduced by the robot part, modular robot structure 
% with small amount of block types, and completely structured discrete environment.
A central process is an act of self-replication -- that is 
an assembly of a replica of itself by the initial robot from a set of parts located on a base.
This process shown in state chart on Figure \ref{fig:LocalAssembly}.

\begin{figure}[h!]
    \centering
    \includegraphics[height=15cm]{LocalAssembly.png} 
    \caption{State chart of self-replication assembly process}
    \label{fig:LocalAssembly}
\end{figure}


Such process can be described as transition systems and is therefore suitable 
for model checking. 
It can be checked for specific requirements, written in form of temporal logic.

In this self-replication process we can derive following actions:
\begin{enumerate}
    \item Connection and disconnection of the end-effector to and from the block or base.
    \item Localization of the part, or receiving of part coordinates from the sensors. 
        This action is a part of "Checking assembly state" and "Searching block in local area"
        in state chart.
    \item Movement of the end-effector from one point to another, with or without connected parts.
\end{enumerate}

Implementation of connection and disconnection actions is platform dependent 
be represented in a control system as boolean state.

Implementation of parts localization actions also depends on specific sensors configuration and  
similarly could be a represented as a boolean array.
Movement action depends mainly on the robot blocks configuration, blocks parameters, 
environmental limitations, initial and target points. 

Model checking of generalized movement, in contrast, has a potential to substantially 
improve the overall self-replication process and adapt it for different 
environment conditions.

In the following sections we will model robot, translate self-replication process 
in transition system and check following properties using NuSMV:
\begin{itemize}
    \item  Check local reachability of some point or a set of points (region) 
        by the robot end-effector. 
    \item Generate a command sequence for a robot to move the end-effector to
        target point, based on counterexamples.
    \item  Research global reachability of some point or region by robot end-effector 
        using the walking form of movements.
    \item Generate a path to a point using walking based on counterexamples.
    \item Analyze set of robot configurations -- check existing configuration or find 
        relevant  configurations of blocks based on counterexamples. 
\end{itemize}

% Analysis of R in R + M + E:
%
% 3. Robot+Environmental reachability checking
%
% Robot configurations synthesis for some objective function
% (defined reachability or for defined environment)
%
% 1. Kinematic robot configuration checking
%
% Search of the optimal robot configuration for defined environment.
% Create a set of configurations combining blocks: 3,4,5,6,7 of different types ()
% Model check reachability of some field (e.g.+-1) for each configuration
%
% Advantage: optimal configuration for a new robot
%
% 2. Model checking of a overall system (control+robot-configuration+
%
% 3. Counterexample based pathfinding in-place assembly
%
% 4. Counterexample based pathfinding in local space
%
% Envitonment:




% plan experiments: 
%     Criteria: success rate, average time, \\
%     1. Test with predefined assembly process (search, move, repeat) 
%     with random distribution of the parts\\
%     2. Test with predefined assembly process (...) with 
%     , criteria and analysis 
%
